\documentclass[10pt]{article}
\usepackage{braket}
\usepackage{soul}
\usepackage{float}
\usepackage{tabularx}
\usepackage{mathtools}
\usepackage{amsmath, amssymb, amsthm}
\usepackage{tikz}
\usetikzlibrary{arrows}
\usetikzlibrary{decorations.text}
\usepackage{enumitem}
\usepackage{geometry}
\usepackage{fancyhdr}
\usepackage{titlesec}
\usepackage{xltabular}
\usepackage{color}
\usepackage{hyperref}
\geometry{margin=1in}

\pagestyle{fancy}
\fancyhf{}
\cfoot{\thepage}
\title{COMPSCI 3AC3 - Assignment 1}
\author{Matthew Farah, \href{mailto:farahm11@mcmaster.ca}{$\langle$farahm11@mcmaster.ca$\rangle$}, 400468251}
\date{\today}

\hypersetup{
	colorlinks=true,
	linkcolor=blue,
	filecolor=magenta,
	urlcolor=blue,
	citecolor=blue,
	anchorcolor=blue
}

\fancyhead{}
\fancyhead[L]{Matthew Farah, \href{mailto:farahm11@mcmaster.ca}{$\langle$farahm11@mcmaster.ca$\rangle$}, 400468251}
\fancyhead[R]{COMPSCI 3AC3 - Assignment 1}

\AtBeginEnvironment{align}{\setcounter{equation}{0}}

\setlength{\parindent}{0cm}

\setcounter{tocdepth}{3}
\hypersetup{bookmarksopen=true, bookmarksopenlevel=2, bookmarksdepth=3}

\newcounter{chapter}
\renewcommand{\thechapter}{\Roman{chapter}}

\newenvironment{chaptertitle}{
	\newpage
	\refstepcounter{chapter}
	\begin{center}
	\vspace*{.5em}
	\par
	\Large
	Part \thechapter
	\phantomsection
	\addcontentsline{toc}{section}{\protect{\large{Part \thechapter}}}
}{
	\end{center}
	\vspace{2.5em}
}

\newcounter{question}
\renewcommand{\thequestion}{\arabic{question}}

\newenvironment{question}{
	\refstepcounter{question}
	\phantomsection
	\addcontentsline{toc}{subsection}{\protect{\large{Question \thequestion}}}
	\vspace{1em}
	\par
	\large\textbf{Question \thequestion}
	\vspace{.2cm}
	\par
	\setcounter{subquestion}{0}
	\setcounter{step}{0}
}{\vspace{.5em}}

\newenvironment{solution}{
	\vspace{1em}\par\large\textbf{{Solution:}}\par\vspace{1em}
}{
	\vspace{1em}
}

\newcounter{subquestion}
\renewcommand{\thesubquestion}{\alph{subquestion}}

\newenvironment{subquestion}{
	\refstepcounter{subquestion}
	\phantomsection
	\vspace{1em}\par\large\textbf{(\thesubquestion)}\hspace{-.5em}
	\setcounter{step}{0}
	\addcontentsline{toc}{subsubsection}{\protect{\large{Part \thequestion.\thesubquestion}}}
}{
	\par
	\vspace{1em}
}

\makeatother

\makeatletter

\newcounter{step}
\renewcommand{\thestep}{\Roman{step}}

\newenvironment{step}{
	\refstepcounter{step}
	\vspace{1.5em}\par\noindent\large\textbf{(\thestep)}
}{
	\par
	\vspace{1.5em}
}

\makeatother

\newcolumntype{Y}{>{\centering\arraybackslash}X}
\renewcommand{\arraystretch}{1.8}

\fontsize{10pt}{14pt}\selectfont

\begin{document}

\maketitle

\tableofcontents
\newpage

\begin{question}
	Show that:
\end{question}

\begin{subquestion}
	$(\alpha + \beta) + \lambda = \alpha + (\beta + \lambda)$ for all
	$\alpha, \beta, \lambda \in \mathbb{C}$.
\end{subquestion}

\begin{solution}

\end{solution}

\begin{subquestion}
	$(\alpha + \beta) + \lambda = \alpha + (\beta + \lambda)$ for all
	$\alpha, \beta, \lambda \in \mathbb{Z} / 2\mathbb{Z}$.
\end{subquestion}

\begin{solution}

\end{solution}

\begin{question}
	Show that:
\end{question}

\begin{subquestion}
	$\lambda(\alpha + \beta) = \lambda\alpha + \lambda\beta$ for all
	$\lambda, \alpha, \beta \in \mathbb{C}$.
\end{subquestion}

\begin{solution}

\end{solution}

\begin{subquestion}
	$\lambda(\alpha + \beta) = \lambda\alpha + \lambda\beta$ for all
	$\lambda, \alpha, \beta \in \mathbb{Z} / 2\mathbb{Z}$.
\end{subquestion}

\begin{solution}

\end{solution}

\begin{question}
	Show that:
\end{question}

\begin{subquestion}
	For every $\alpha \in \mathbb{C}$ with $\alpha \neq 0$, there exists a
	unique $\beta \in \mathbb{C}$ such that $\alpha\beta = 1$.
\end{subquestion}

\begin{solution}

\end{solution}

\begin{subquestion}
	For every $\alpha \in \mathbb{Z} / 2\mathbb{Z}$ with
	$\alpha \neq 0$, there exists a unique $\beta \in \mathbb{C}$ such that
	$\alpha\beta = 1$.
\end{subquestion}

\begin{solution}

\end{solution}

\begin{question}
	Show that $(ab)x = a(bx)$ for all $x \in \mathbb{F}^n$, where
	$\mathbb{F}$ is any field.
\end{question}

\begin{solution}

\end{solution}

\begin{question}
	Show that $(\alpha + \beta)x = \alpha{x} + \beta{x}$ for all $a, b \in
	\mathbb{F}$ and all $x \in \mathbb{F}^n$, where $\mathbb{F}$ is any
	field.
\end{question}

\begin{solution}

\end{solution}

\end{document}
